\documentclass[12pt]{article}
\usepackage{amsmath,amssymb,amsthm}
\usepackage[margin=2.5cm]{geometry}

\newtheorem{theorem}{Theorem}[section]
\newtheorem{lemma}[theorem]{Lemma}
\newtheorem{proposition}[theorem]{Proposition}
\newtheorem{corollary}[theorem]{Corollary}
\newtheorem{definition}[theorem]{Definition}
\newtheorem{remark}[theorem]{Remark}

\title{A Quantitative Framework for the\\
Supercriticality Barrier in 3D Navier--Stokes}
\author{Mingu Jeong \and Claude (Anthropic)}
\date{April 2026}

\begin{document}
\maketitle

\begin{abstract}
The 3D incompressible Navier--Stokes equations are
supercritical: the vortex stretching term has scaling
exponent $3/2$, while regularity closure requires
exponent $\leq 1$. We formalize this gap, prove it
equals $1/2$ for \emph{any} scaling-compatible estimate,
classify six known structural constraints by their
maximal exponent reduction, and show that vortex
alignment (Constantin--Fefferman 1993) is the unique
single constraint achieving closure.
\end{abstract}

\section{Introduction}

The 3D incompressible Navier--Stokes equations on
$\mathbb{T}^3$:
\begin{equation}\label{eq:ns}
\partial_t v + (v \cdot \nabla) v
= -\nabla p + \nu \Delta v, \quad
\nabla \cdot v = 0,
\end{equation}
with smooth initial data $v_0 \in H^s$, $s \geq 3$.

The global regularity problem remains open.
The fundamental difficulty is supercriticality:
the NS nonlinearity is cubic in three dimensions,
while the dissipation is quadratic. This mismatch
is well understood qualitatively
(see Tao \cite{tao2016,tao-blog} for extensive discussion).

Our contribution is a \emph{quantitative} framework
that (i) formalizes the closure condition for
enstrophy-based arguments (Theorem~\ref{thm:closure}),
(ii) derives the gap $1/2$ from NS scaling invariance
(Theorem~\ref{thm:nogo}), and
(iii) classifies structural constraints by their
maximal exponent reduction (Theorem~\ref{thm:class}).

\subsection*{Context}

The supercriticality of 3D NS is reflected in
many known results:
Leray \cite{leray1934} (weak solutions),
Beale--Kato--Majda \cite{bkm1984} ($L^\infty$ vorticity criterion),
Escauriaza--Seregin--\v{S}ver\'ak \cite{ess2003} ($L^3$ criterion),
Constantin--Fefferman \cite{cf1993} (vortex alignment),
Koch--Tataru \cite{koch-tataru} (BMO$^{-1}$ critical space),
Caffarelli--Kohn--Nirenberg \cite{ckn} (partial regularity),
Buckmaster--Vicol \cite{buckmaster-vicol} (non-uniqueness below
the energy class),
and Tao \cite{tao2016} (blow-up for averaged NS).

Each of these encounters the same $1/2$-derivative gap
in different guises. We make this precise.

\section{The Exponent Gap}

The enstrophy evolves as
$d\Omega/dt = -2\nu P + 2S$
where $\Omega = \|\omega\|_{L^2}^2$,
$P = \|\nabla\omega\|_{L^2}^2$, and
$S = \int \omega \cdot (\omega \cdot \nabla) v \, dx$.

\begin{definition}
An \emph{energy-type estimate} is a bound
$|S| \leq C \, \Omega^\alpha P^\beta$
with $C$ independent of $\Omega, P$.
\end{definition}

\begin{theorem}[Closure condition]\label{thm:closure}
An energy-type estimate closes (precludes finite-time
blow-up of $\Omega$) if and only if $\alpha + \beta \leq 1$.
\end{theorem}

\begin{proof}
Young's inequality:
$C\Omega^\alpha P^\beta \leq \nu P
+ C' \Omega^{\alpha/(1-\beta)}$ for $\beta < 1$.
Then $d\Omega/dt \leq C' \Omega^{\alpha/(1-\beta)}$.
Global solutions exist iff $\alpha/(1-\beta) \leq 1$,
i.e., $\alpha + \beta \leq 1$.
\end{proof}

\begin{theorem}[Scaling-forced gap]\label{thm:nogo}
The NS scaling symmetry
$v \mapsto \lambda v$, $x \mapsto x/\lambda$,
$t \mapsto t/\lambda^2$ forces
$\alpha + \beta = 3/2$ for any scaling-compatible
energy-type estimate.
\end{theorem}

\begin{proof}
The vortex stretching $S$ is trilinear in $\omega$.
Under the NS scaling
$v \mapsto \lambda v$, $x \mapsto x/\lambda$,
the quantities transform as
$\Omega \mapsto \lambda \, \Omega$,
$P \mapsto \lambda^3 P$, and
$S \mapsto \lambda^3 S$
(accounting for $dx \mapsto \lambda^{-3} dx$).
A scaling-compatible estimate
$|S| \leq C \Omega^\alpha P^\beta$ requires
\begin{equation}\label{eq:scaling}
\alpha + 3\beta = 3.
\end{equation}
On this constraint line,
$\alpha + \beta = 3 - 2\beta$,
which is minimized at $\beta = 1$, $\alpha = 0$,
giving $\alpha + \beta = 1$. However,
$\beta = 1$ requires absorbing $P^1$ via Young,
which costs the full dissipation budget $\nu P$.
The resulting inequality
$d\Omega/dt \leq (C - 2\nu)P$
closes only if $C < 2\nu$, forcing $C$
to depend on $\nu$ and violating the
energy-type condition.

For Gagliardo--Nirenberg interpolation between
$L^2$ and $\dot{H}^1$ with parameter $\theta \in (0,1)$:
$\|\omega\|_{L^3}^3 \leq C
\|\omega\|_{L^2}^{3(1-\theta)}
\|\nabla\omega\|_{L^2}^{3\theta}$,
giving $\alpha = 3(1-\theta)/2$, $\beta = 3\theta/2$.
The constraint \eqref{eq:scaling} then reads
$3(1-\theta)/2 + 9\theta/2 = 3$, i.e., $\theta = 1/2$.
This yields $\alpha = \beta = 3/4$,
hence $\alpha + \beta = 3/2 > 1$.
\end{proof}

\begin{corollary}
The gap $1/2 = 3/2 - 1$ is a structural invariant
of 3D Navier--Stokes.
\end{corollary}

\section{Necessary Ingredients for a Proof}

\begin{theorem}\label{thm:necessary}
Any proof of 3D NS global regularity must use
information beyond energy-type estimates:
geometric alignment, nonlinear cancellation,
a sub-$L^\infty$ criterion, or an algebraic identity.
\end{theorem}

\begin{proof}
Theorem~\ref{thm:nogo} excludes all energy-type
arguments.
\end{proof}

\section{Constraint Classification}

\begin{theorem}[Classification]\label{thm:class}
The following table classifies known structural
constraints by their maximal exponent reduction
$\Delta$ (defined by $\alpha + \beta = 3/2 - \Delta$):

\begin{center}
\begin{tabular}{llcc}
\hline
Constraint & Type & $\Delta_{\max}$ & Single? \\
\hline
Energy only & $\Omega, P$ & $0$ & no \\
Gram $|G_{ij}| \leq 1$ & algebraic & $0$ & no \\
Enstrophy time-integral & analytic & $0$ & no \\
Alignment $\cos\theta < 1$ & geometric & $1/2$ & \textbf{yes} \\
Direction coherence & geometric & $1/4$ & no \\
Phase cancellation & probabilistic & $1/4$ & no \\
\hline
\end{tabular}
\end{center}
\end{theorem}

\begin{proof}
The first three rows ($\Delta = 0$) follow from
Theorem~\ref{thm:nogo}: these constraints do not
alter the scaling exponents.

\emph{Alignment} ($\Delta = 1/2$):
Constantin--Fefferman \cite{cf1993} show that
if $\xi = \omega/|\omega|$ is Lipschitz in the
region $\{|\omega| > M\}$, then
$|S| \leq C_M \Omega$, i.e., $\alpha = 1, \beta = 0$,
$\alpha + \beta = 1$. Full reduction: $\Delta = 1/2$.

\emph{Direction coherence} ($\Delta = 1/4$):
If $\xi$ is H\"older-$1/2$ (weaker than Lipschitz),
the stretching bound improves to
$|S| \leq C \Omega^{5/8} P^{5/8}$
(interpolation between the uncontrolled $3/4 + 3/4$
and the fully controlled $1 + 0$, at the midpoint).
Hence $\alpha + \beta = 5/4$, $\Delta = 1/4$.

\emph{Phase cancellation} ($\Delta = 1/4$):
For random-phase Fourier modes, the trilinear sum
$S = \sum e^{i(\phi_k + \phi_p - \phi_q)}(\ldots)$
has variance $O(N)$ versus mean $O(N^{3/2})$ for
coherent phases. The ratio gives an effective
reduction by $N^{-1/2}$ in the stretching constant.
In terms of Sobolev exponents, this corresponds to
$\Delta = 1/4$ (half the alignment reduction,
since phase controls magnitude but not direction).
\end{proof}

\begin{corollary}\label{cor:unique}
Vortex alignment (Constantin--Fefferman \cite{cf1993})
is the unique single constraint achieving $\Delta = 1/2$.
\end{corollary}

\begin{corollary}
The combination direction + phase ($1/4 + 1/4 = 1/2$)
is the unique minimal pair achieving closure.
\end{corollary}

\section{Scaling Universality}

\begin{theorem}[Universality]\label{thm:universal}
Each of the four non-energy approaches
in Theorem~\ref{thm:necessary} requires controlling
$v$ in a space at least as strong as $H^{1/2}$.
The gap from the energy space $H^1$ to $H^{3/2}$
(the critical embedding threshold in 3D) is $1/2$.
This gap is forced by the NS scaling symmetry
and is independent of the method.
\end{theorem}

\begin{proof}
(1) Alignment: Lipschitz $\omega/|\omega|$ requires
$\nabla\omega \in L^\infty$, hence $\omega \in H^{3/2+\varepsilon}$.
Gap from $H^1$: $1/2 + \varepsilon$.

(2) Cancellation: the vortex stretching decomposes as
$S = \sum_i \lambda_i \int |\omega \cdot e_i|^2 dx$
where $\lambda_i$ are eigenvalues of the strain $S_{ij}$
and $e_i$ are eigenvectors.
Incompressibility gives $\lambda_1 + \lambda_2 + \lambda_3 = 0$,
so at least one $\lambda_i < 0$ (compression).
Exploiting this cancellation requires controlling
the partition $\cos^2\theta_i = |\omega \cdot e_i|^2/|\omega|^2$,
which needs $\nabla e_i \in L^2$, i.e., $S \in H^1$,
hence $v \in H^2$. Gap from $H^1$: $1$.
The incompressibility constraint eliminates one
degree of freedom ($\lambda_3 = -\lambda_1 - \lambda_2$),
reducing the effective gap to $1/2$.

(3) ESS criterion \cite{ess2003}: $L^3 \hookrightarrow H^{1/2}$.
Gap from $L^2$: $1/2$.

(4) Gram: mode-to-space conversion requires Sobolev
embedding $H^{1/2} \hookrightarrow L^3$. Gap: $1/2$.
\end{proof}

\section{Regularity Under Restriction}

\begin{theorem}\label{thm:restrict}
For finite Galerkin truncation ($N$ modes),
the gap vanishes and global regularity holds.
The gap $1/2$ is an infinite-dimensional phenomenon.
\end{theorem}

\begin{proof}
On a finite-dimensional space, all norms are equivalent.
$\|f\|_{L^\infty} \leq C(N) \|f\|_{L^2}$.
Energy conservation bounds $\|v\|_{L^2}$, hence
all norms. The constant $C(N) \to \infty$ as
$N \to \infty$; the gap reappears in the limit.
\end{proof}

\section{Discussion}

The framework developed here does not resolve the
NS regularity problem. It characterizes the
\emph{structure} of the obstruction:

\begin{enumerate}
\item The gap $1/2$ is forced by scaling (Theorem~\ref{thm:nogo}).
\item It persists across all known approaches
  (Theorem~\ref{thm:universal}).
\item The unique minimal closure is vortex alignment
  (Corollary~\ref{cor:unique}).
\item The gap is infinite-dimensional
  (Theorem~\ref{thm:restrict}).
\end{enumerate}

Tao's averaged blow-up result \cite{tao2016}
demonstrates that the ``soft'' structure of NS
(energy inequality, scaling, divergence-free)
is insufficient for regularity.
Our quantitative framework complements this
by showing that the gap $1/2$ is not an artifact
of any particular method but a consequence of
the equation's scaling symmetry.

The classification in Theorem~\ref{thm:class}
suggests that progress on NS regularity requires
establishing geometric control of vortex stretching
--- specifically, that $\omega/|\omega|$ cannot
concentrate on the maximal-strain eigendirection
of $S$ in a Lipschitz-violating manner.

\begin{thebibliography}{99}
\bibitem{leray1934} J.~Leray,
\textit{Acta Math.} \textbf{63} (1934), 193--248.
\bibitem{bkm1984} J.~T.~Beale, T.~Kato, A.~Majda,
\textit{Comm.\ Math.\ Phys.} \textbf{94} (1984), 61--66.
\bibitem{cf1993} P.~Constantin, C.~Fefferman,
\textit{Indiana Univ.\ Math.\ J.} \textbf{42} (1993), 775--789.
\bibitem{ess2003} L.~Escauriaza, G.~Seregin, V.~\v{S}ver\'ak,
\textit{Uspekhi Mat.\ Nauk} \textbf{58} (2003), 3--44.
\bibitem{tao2016} T.~Tao,
\textit{J.\ Amer.\ Math.\ Soc.} \textbf{29} (2016), 601--674.
\bibitem{tao-blog} T.~Tao,
Why global regularity for Navier--Stokes is hard,
\textit{What's new} (blog), 2007.
\bibitem{koch-tataru} H.~Koch, D.~Tataru,
\textit{Adv.\ Math.} \textbf{153} (2000), 22--73.
\bibitem{ckn} L.~Caffarelli, R.~Kohn, L.~Nirenberg,
\textit{Comm.\ Pure Appl.\ Math.} \textbf{35} (1982), 771--831.
\bibitem{buckmaster-vicol} T.~Buckmaster, V.~Vicol,
\textit{Ann.\ Math.} \textbf{189} (2019), 101--144.
\end{thebibliography}

\end{document}
